\section{战略转型、融资历史与潜在融资需求}
\label{sec:financing}

本章回答第二部分最后、也是最实际的一个问题：\hl{在周期底部，
这 \num{104} 家公司各自需要什么，又能从哪里得到它}？
分析方法是从两类可观测的行为数据出发——\emph{公告}（公司自己披露的意图）
与\emph{财务结构}（公司实际承受的约束），
再把两者交叉，得到融资需求的三个层次。

\subsection{公告中的行为信号}
\label{sec:financing-signals}

近三年（2023 年 9 月---2026 年 9 月）\num{104} 家公司累计披露公告 \num{41006} 条。
按标题关键词分类后得到的事件分布如下：

\begin{table}[htbp]
\centering
\footnotesize
\caption{农业上市公司近三年公告的事件分类统计（按标题关键词归类，共 \num{41006} 条；同一公告可命中多个类别，各行数值不可相加）}
\label{tab:fin-events}
\begin{tabular}{@{}p{2.9cm}rp{7.3cm}@{}}
\toprule
事件类别 & 公告数 & 典型标题特征与解读 \\
\midrule
担保与关联交易 & \num{3184} & “为子公司提供担保”“关联交易”——集团内输血，隐含资金压力 \\
回购与激励 & \num{1853} & “回购股份”“股权激励”“员工持股”——管理层对估值与留人的判断 \\
再融资 & \num{1363} & “向特定对象发行”“募集说明书”“募集资金存放与使用”——融资动作 \\
股东行为 & \num{1353} & “减持”“增持”“质押”——股东层面的流动性需求 \\
业绩与风险 & \num{1252} & “业绩预告”“风险提示”“立案/问询”——经营与合规风险披露 \\
重大重组与并购 & \num{348} & “重大资产重组”“收购”“股权转让”——外延扩张或保壳 \\
对外投资与转型 & \num{274} & “设立子公司”“战略合作”“框架协议”“扩产”——产能与赛道调整 \\
\bottomrule
\end{tabular}
\end{table}

\hl{“担保与关联交易”居首（\num{3184} 条）具有明确的经济含义}：
在农业集团企业内部，母公司为亏损子公司提供担保是维持其
银行授信与生产经营的常规手段。担保公告越密集，通常意味着\emph{集团内部
对单一板块（多为生猪养殖）的资金依赖度越高}。
“再融资”类公告 \num{1363} 条则直接对应资本市场融资动作，
其中 \textbf{新希望}（\num{95} 条）、\textbf{国投中鲁}（\num{85} 条）、
\textbf{*ST广糖}（\num{78} 条）、\textbf{巨星农牧}（\num{75} 条）居前四位，
\hl{这四家公司在 2026 年上半年均处于亏损或微利状态}
（归母净利润分别为 \num{-17.02}、\num{0.02}、\num{-0.16} 与 \num{-8.38} 亿元）。

按三级行业看，公告密度最高的是\textbf{生猪养殖}（公司均再融资公告 \num{28.3} 条、
中位数 \num{27} 条，居全部三级行业首位），其次是\textbf{果蔬加工}（\num{27.6} 条）；
若看重组与并购类公告，密度最高的是\textbf{果蔬加工}（\num{10.6} 条）与\textbf{林业}（\num{9.0} 条）。
前者的高密度对应全行业性的补血需求，后两者则集中在\hl{重组与壳资源运作}上——
国投中鲁、\textbf{*ST景谷}、中基健康、永安林业均属此类，
从净资产与公告结构看，其共同特征是主业萎缩、壳价值可能高于经营价值。

\subsection{融资历史：股本与分红的两个事实}
\label{sec:financing-history}

\begin{itemize}[leftmargin=1.4em, itemsep=2pt]
  \item \textbf{股本变动是普遍事实}：\num{104} 家公司\hl{全部在近三年存在股本变更记录}
        （含送转、股权激励行权、定增、可转债转股）——\emph{没有任何一家的股本在三年内保持不变}。
  \item \textbf{分红能力分化}：\num{73} 家公司在近三年实施过现金分红，
        \hl{\num{31} 家连续三年未分红}（占 \num{29.8}\%），
        其中多数同时处于亏损状态。
\end{itemize}

这两个事实合起来说明一个规律：\hl{农业上市公司的股本变动频率
远高于其分红频率}。在周期底部，公司倾向于\emph{用股权而不是现金流来解决问题}——
定增、可转债、股权激励行权都会增加股本，而现金分红是唯一
真正返还股东的渠道，因此在下行周期中最先被削减。
\emph{这构成了“未来潜在融资需求”分析的基本背景}：
融资的动机在周期底部是刚性（补血），而不是可选（扩张）。

\subsection{重大战略转型的三条路径}
\label{sec:financing-strategy}

从公告与财务数据的组合中，可以归纳出本轮周期中农业上市公司转型的三条路径。

\paragraph{路径一：一体化收缩（“把亏损的业务砍掉”）。}
以\textbf{新希望}为代表：公司以饲料起家、在 2020---2021 年猪价高点
举债扩张生猪产能，2023---2025 年连续三年将生猪板块作为亏损主体，
2026 年上半年仍亏损 \num{17.02} 亿元、2026 年 6 月末资产负债率 \num{70.7}\%、
再融资相关公告 \num{95} 条（全样本第一）。
其战略实质是\hl{“以融资换时间”}：通过定增、出售资产、引入战投
维持资产负债表，等待猪价回升。
同类公司包括\textbf{唐人神}（\num{-10.46} 亿元）、
\textbf{大北农}（\num{-6.77} 亿元）、\textbf{天邦食品}（\num{-15.43} 亿元）。

\paragraph{路径二：跨界与保壳（“换一个行业重新开始”）。}
以\textbf{国投中鲁}（\num{85} 条再融资公告 + \num{32} 条重组公告，全样本最高重组密度）、
\textbf{绿康生化}（\num{23} 条重组公告，从动保转向微生物发酵与原料药）、
\textbf{*ST景谷}（\num{23} 条重组公告）为代表。
这类公司的特征是：\hl{主业规模小、资产轻、市值仍有一定支撑}，
转型方向多为与农业无关的新兴产业。
\emph{它们的融资需求是“保壳型”的}——不是为了扩产，而是为了维持上市地位。

\paragraph{路径三：向消费端与高附加值环节延伸（“把周期变成品牌”）。}
以\textbf{海大集团}（服务型饲料，半年 ROE \num{6.5}\%、未年化）、
\textbf{乖宝宠物}（半年 ROE \num{4.1}\%、半年营收同比 \num{+10.0}\%）、
\textbf{中宠股份}（半年 ROE \num{4.0}\%、半年营收同比 \num{+34.9}\%）、
\textbf{瑞普生物}（动保平台化，半年 ROE \num{4.3}\%）
为代表。它们的共同点是\hl{把收入结构从“价格 $\times$ 数量”
转向“品牌 $\times$ 复购”}，从而在农产品价格下行时保持利润率。
这类公司是本轮周期中唯一\emph{以扩张为目的}融资的群体——
乖宝宠物（2023 年 8 月）、索宝蛋白（2023 年 12 月）
是全部 \num{104} 家公司中上市最晚的两家。

\subsection{未来潜在融资需求的分层}
\label{sec:financing-need}

把财务约束与转型路径交叉，可以把潜在融资需求分为三档。

\begin{figure}[htbp]
\centering
\begin{tikzpicture}
\begin{groupplot}[
  group style={group size=2 by 1, horizontal sep=1.6cm},
  agri axis, width=5.5cm, height=5.4cm,
]
\nextgroupplot[
  xlabel={2026 年 6 月末资产负债率（\%）}, ylabel={2026 年上半年 ROE（\%）},
  xlabel style={font=\scriptsize}, ylabel style={font=\scriptsize},
  tick label style={font=\scriptsize},
  % 允许超出坐标范围：天邦食品的半年 ROE（-96.8%）已在图注中说明，
  % 若纳入坐标轴会把其余公司的主体分布压缩到不足图高的十分之一
  xmin=0, xmax=105, ymin=-70, ymax=35,
  grid=both, grid style={InkGray!18},
  title style={font=\small}, title={负债率与回报的负相关},
]
\addplot[only marks, mark=*, mark size=1.9pt, color=OilBlue, opacity=0.85]
  table[x=debt, y=roe]{data/clean/comp_debt_roe.dat};
\draw[SoftRed, dashed, line width=0.6pt] (axis cs:65,-70) -- (axis cs:65,35);
\node[font=\tiny, SoftRed, rotate=90, anchor=south] at (axis cs:63,-20) {65\% 警戒线};
\nextgroupplot[
  xbar, xmin=0, xmax=105, ymin=-0.7, ymax=13.7,
  ytick={0,1,2,3,4,5,6,7,8,9,10,11,12,13},
  yticklabels from table={data/clean/comp_risk.dat}{label},
  yticklabel style={font=\tiny, align=right, text width=1.25cm},
  xlabel={2026 年 6 月末资产负债率（\%）}, xlabel style={font=\scriptsize},
  tick label style={font=\scriptsize},
  bar width=5.5pt,
  title style={font=\small}, title={亏损且负债率$>$65\%的公司},
]
\addplot[fill=SoftRed!70, draw=SoftRed] table[y=i, x=debt]{data/clean/comp_risk.dat};
\end{groupplot}
\end{tikzpicture}
\caption{农业上市公司的财务约束：资产负债率与 ROE 的关系（左，2026 年半年报口径）与高负债亏损公司（右；左图 \num{103} 家可算 ROE，其中 \num{1} 家（天邦食品，\num{-96.8}\%）超出纵轴范围未显示）}
\label{fig:fin-risk}
\end{figure}

左图给出一条清晰的经验规律：\hl{资产负债率超过 \num{65}\% 后，
ROE 几乎全部落在负区间}；\num{103} 家可算 ROE 的公司中，
半年 ROE 超过 \num{5}\%（约相当于年化 \num{10}\%）的有 \num{16} 家，
其中 \num{10} 家的资产负债率低于 \num{65}\%——\emph{高负债与高回报在样本中几乎互斥}。
右图则列出同时满足“2026 年上半年亏损 + 6 月末负债率 $>$ \num{65}\%”的 \num{14} 家公司，
其中 \hl{\num{6} 家负债率超过 \num{80}\%}：
中基健康（\num{104.1}\%）、*ST广糖（\num{100.2}\%）、天邦食品（\num{91.2}\%）、
新五丰（\num{84.2}\%）、金新农（\num{81.5}\%）、天域生物（\num{81.0}\%）。
\emph{中基健康与*ST广糖的净资产已经为负，其余四家的负债率也已接近
“净资产被亏损侵蚀殆尽”的临界点}。

据此，未来潜在融资需求可分为三档：

\begin{table}[htbp]
\centering
\footnotesize
\caption{农业上市公司未来潜在融资需求的三层结构（基于 2026 年半年报与近三年公告数据）}
\label{tab:fin-layers}
\begin{tabular}{@{}p{1.7cm}p{3.3cm}p{4.6cm}p{4.6cm}@{}}
\toprule
层次 & 特征与公司数 & 融资动机 & 可行的融资方式 \\
\midrule
第一档\newline 补血型 & 2026 年上半年亏损且 6 月末负债率 $>$65\%：\newline \num{14} 家（其中 $>$80\% \num{6} 家）
  & 维持银行授信、偿还到期债务、避免退市
  & 定增（需股价配合）、债务重组、出售资产、引入国资或战投；\hl{股权融资是最现实选择} \\
第二档\newline 保壳型 & 重组并购公告 $\geq$ \num{10} 条：\newline \num{9} 家（其中 2026 上半年亏损 \num{6} 家）
  & 资产置换、跨界注入新业务、维持上市地位
  & 重大资产重组、发行股份购买资产、控制权转让；\hl{融资与“易主”同时发生} \\
第三档\newline 扩张型 & 半年 ROE $>$ \num{5}\%（约合年化 \num{10}\%）且 6 月末负债率 $<$ \num{65}\%：\newline \num{10} 家（动保、细分加工与农资为主）
  & 产能扩张、海外布局、品类并购
  & 定增 + 可转债、并购贷款；\hl{这类融资会被市场定价为成长而非补血} \\
\bottomrule
\end{tabular}
\end{table}

需要说明的是三档之间的\emph{迁移}：与 2025 年年报口径相比，第一档的公司数由 \num{18} 家
减少到 \num{14} 家（其中 \num{6} 家负债率超过 \num{80}\%），
原因是个别高负债公司在 2026 年上半年实现扭亏或完成了资产处置，
但全样本亏损面却由 \num{34} 家扩大到 \num{46} 家（见第 \ref{sec:comp-h1} 节）——
\hl{“头部补血、尾部出清”是本轮周期底部最典型的资产负债表形态}。

按金额量级估算，第一档的融资需求最为迫切：这 \num{14} 家公司的
2026 年上半年亏损合计约 \hl{\num{73.1} 亿元}
（此处为该档成员的亏损合计数，并非板块或全样本的亏损合计）。
若只按“补足一年经营亏损 + 覆盖一年利息支出”的最低口径，
年度资金缺口约在 \hl{\num{150}---\num{250} 亿元}量级；
若按“把资产负债率压回 \num{65}\% 警戒线”的重组口径，
这 \num{14} 家公司 2026 年 6 月末的负债合计为 \num{1855} 亿元、
净资产合计仅 \num{694} 亿元，\hl{需要补充的净资产将超过 \num{2000} 亿元}
（按 $\text{目标净资产} = \text{负债} \div 0.65$ 粗略推算）。
而第三档公司的融资规模取决于其扩张节奏，单笔通常在
\num{5}---\num{20} 亿元之间。\emph{两者的根本差别在于：
第一档融资的用途是“活下去”，第三档融资的用途是“长大”}。

\begin{keybox}[融资需求分析的三个结论]
\normalsize
\textbf{① 融资动机由周期位置决定，而非行业属性}。同一家公司在 2021 年
募资用于扩产，在 2026 年募资用于补血——\hl{公告主题的变化就是周期的可见痕迹}。\par\vspace{3pt}
\textbf{② 高负债与亏损高度重合}：\num{46} 家亏损公司中有 \num{14} 家同时负债率超过 \num{65}\%，
其中 \num{6} 家超过 \num{80}\%。
\hl{负债率是比净利润更早的预警指标}。\par\vspace{3pt}
\textbf{③ 潜在融资需求的规模在数百亿元量级}，
且集中在生猪产业链（新希望、天邦食品、新五丰、唐人神、巨星农牧）
与困境加工企业（*ST广糖、中基健康、大禹节水、百洋股份）。
\hl{资本市场对农业板块的融资功能，在未来两三年内将主要被用于修复资产负债表，
而不是扩大产能}。
\end{keybox}
